\exercisesection{2}

\begin{enumerate}
\def\labelenumi{\arabic{enumi}.}
\item
  Give an example of an injective ring homomorphism \(\rho: \mathbf{A} \to \mathbf{B}\) such that \(\mathbf{B}\) is finite A-algebra, the mapping \(^a\rho: \operatorname{Spec}(\mathbf{B}) \to \operatorname{Spec}(\mathbf{A})\) is surjective but \(\mathbf{B}\) is not a flat A-module. Conversely, give an example of a faithfully flat A-algebra which is finitely generated but not integral over \(\mathbf{A}\) .
\item
  Let A be a ring such that \(\text{Spec}(A)\) is Hausdorff (Chapter II, § 4, Exercise 16). Show that, for every A-algebra B integral over A, \(\text{Spec}(B)\) is Hausdorff.
\item
  Let A be a ring and A' a ring containing A and integral over A. For every ideal a of A, show that the radical of a is the intersection of A and the radical of aA'. Show that the latter is the set of elements \(x \in A'\) satisfying an equation of integral dependence whose coefficients (other than the dominant coefficient) belong to a (note that, if y, z are two elements of A' satisfying such an equation of integral dependence, so does y - z).
\item
  Let A be a Noetherian ring, \(\rho: A \to B\) a ring homomorphism making B into a finitely generated A-module and N a finitely generated B-module. Show that the prime ideals \(\mathfrak{p} \in \operatorname{Ass}(\rho_{*}(N))\) are the inverse images under \(\rho\) of the prime ideals \(\mathfrak{q} \in \operatorname{Ass}(N)\) .
\item
  Let K be a field; in the polynomial domain K{[}X, Y{]} in two indeterminates consider the subrings A = K{[}X \(^{4}\) , Y \(^{4}\) {]}, A' = K{[}X \(^{4}\) , X \(^{3}\) Y, XY \(^{3}\) , Y \(^{4}\) {]}; A is Noetherian and A' is a finite A-algebra. Let p = AY \(^{4}\) , which is a prime ideal of A; show that m = A'X \(^{4}\) + A'X \(^{3}\) Y + A'XY \(^{3}\) + A'Y \(^{4}\) is an (immersed) prime ideal associated with the ideal pA' of A' but does not lie above p (note that m is the annihilator of the class of X \(^{2}\) Y \(^{6}\) in the ring A'/pA') (cf.~§ 1, Exercise 26(d)).
\item
  Define an increasing sequence of algebraic number fields \(K_{0} = Q, K_{1}, \ldots, K_{n}, \ldots\) such that, if \(A_{n}\) denotes the integral closure of Z in \(K_{n}, p\) a prime number \(\neq 2\) , \(K_{n}\) is a quadratic extension of \(K_{n-1}\) and there exists in \(A_{n}\) a prime ideal \(p_{n}\) lying above (p) such that there are in \(A_{n+1}\) two distinct prime ideals \(p_{n+1}, p'_{n+1}\) lying above \(p_{n}\) (observe that, in a finite field of characteristic \(\neq 2\) , there are always elements which are not squares). Let K be the union of the \(K_{n}\) and \(A'\) the integral closure of Z in K; show that there exists an infinity of (maximal) ideals of \(A'\) lying above the maximal ideal (p) of Z.
\item
  Let A be a Noetherian integral domain and A' its integral closure. Show that for every prime ideal p of A there exists only a finite number of prime ideals of A' lying above p.~(Reduce it first to the case where A is a local ring with maximal ideal p; consider its completion \(\hat{A}\) and the total ring of fractions B of \(\hat{A}\) and note that the field of fractions K of A is identified with a subring of B (Chapter III, § 3, no. 4, Corollary 2 to Theorem 3). Note that the spectrum of B is finite and therefore that in B there is only a finite number of idempotents which are orthogonal to one another. Then let C \(\supset A\) be a subring of K which is a finitely generated A-module and therefore semi-local. Note that the completion \(\hat{C}\) of C is identified with a subring of B; conclude by noting that \(\hat{C}\) is the product of a finite number of local rings, this number being equal to the number of maximal ideals of C). Generalize to the integral closure of a Noetherian ring in its total ring of fractions (reduce it to the case of a reduced ring).
\item
  A local ring A is called unibranch if its integral closure is a local ring.
\end{enumerate}

\begin{enumerate}
\def\labelenumi{(\alph{enumi})}
\item
  Let A be a local integral domain and K its field of fractions. For A to be unibranch, it is necessary and sufficient that every sub-A-algebra of K which is a finitely generated A-module be a local ring.
\item
  Show that, if the completion of a Noetherian local ring A is an integral domain, A is unibranch (if B is a finite A-algebra contained in K, show that the subring of the field of fractions of \(\hat{A}\) generated by B and \(\hat{A}\) is isomorphic to \(\mathbf{B} \otimes_{\mathbf{A}} \hat{\mathbf{A}}\) , using the flatness of the A-module \(\hat{\mathbf{A}}\) and the fact that \(\mathbf{B}\) is contained in a finitely generated free A-module contained in \(\mathbf{K}\) ; then use (a) and Chapter III, § 2, no. 13, Corollary to Proposition 19).
\end{enumerate}

\begin{enumerate}
\def\labelenumi{\arabic{enumi}.}
\setcounter{enumi}{8}
\tightlist
\item
  Let \(k_0\) be a field and \(A'\) the polynomial ring \(k_0[X_n]_{n \in \mathbb{N}}\) in an infinite sequence of indeterminates; it is an integrally closed domain whose field of fractions is \(K' = k_0(X_n)_{n \in \mathbb{N}}\) . Let \(\sigma\) denote the \(k_0\) -automorphism of \(A'\) such that
\end{enumerate}

\[
\sigma (\mathrm{X} _ {2 n}) = - \mathrm{X} _ {2 n} + \mathrm{X} _ {2 n + 1}, \quad \sigma (\mathrm{X} _ {2 n + 1}) = \mathrm{X} _ {2 n + 1}
\]

for all \(n \geqslant 0\) ; \(\sigma\) and the identity automorphism form a group G of automorphisms of A'. Let A be the ring of invariants of G and let K be its field of fractions; K' is a separable extension of K of degree 2 and A' is the integral closure of A in K'.

\begin{enumerate}
\def\labelenumi{(\alph{enumi})}
\tightlist
\item
  Show that \(\mathbf{A}'\) is an A-module admitting an infinite basis. (If \(k_0\) is of characteristic \(\neq 2\) , show that, if we write \(\mathbf{Y}_n = \mathbf{X}_{2n} - \frac{1}{2}\mathbf{X}_{2n+1}\) , \(\mathbf{A}\) is the subring of \(\mathbf{A}'\) generated by the \(\mathbf{X}_{2n+1}\) and the products \(\mathbf{Y}_n\mathbf{Y}_m\) for \(n \geqslant 0\) , \(m \geqslant 0\) . If \(k_0\) is of characteristic 2, \(\mathbf{A}'\) is generated by the \(\mathbf{X}_{2n+1}\) and the
\end{enumerate}

\[
\mathrm{X} _ {2 n} ^ {2} + \mathrm{X} _ {2 n} \mathrm{X} _ {2 n + 1}
\]

for \(n\geqslant 0.\)

\begin{enumerate}
\def\labelenumi{(\alph{enumi})}
\setcounter{enumi}{1}
\item
  Let \(\mathfrak{m}'\) be the maximal ideal of \(\mathbf{A}'\) generated by all the \(\mathbf{X}_n\) and let \(\mathfrak{m} = \mathbf{A} \cap \mathfrak{m}'\) ; show that \(\mathfrak{m}' = \mathfrak{mA}'\) , that \(\mathfrak{m}'\) is the only ideal of \(\mathbf{A}'\) lying above \(\mathfrak{m}\) and that the fields \(\mathbf{A} / \mathfrak{m}\) and \(\mathbf{A}' / \mathfrak{m}'\) are canonically isomorphic.
\item
  Let \(\mathfrak{p}'\) be the prime ideal of \(\mathbf{A}'\) generated by the \(\mathbf{X}_{2n+1}\) for \(n \geqslant 0\) and let \(\mathfrak{p} = \mathfrak{p}' \cap \mathbf{A}\) ; let \(k, k'\) be the fields of fractions of \(\mathbf{A}/\mathfrak{p}\) and \(\mathbf{A}'/\mathfrak{p}'\) respectively; \(\mathfrak{p}'\) is the only prime ideal of \(\mathbf{A}'\) lying above \(\mathfrak{p}\) . If \(k_0\) is not of characteristic 2, \(k'\) is a separable extension of \(k\) of degree 2; if \(k_0\) is of degree 2, \(k'\) is an infinite radical extension of \(k\) .
\end{enumerate}

\begin{enumerate}
\def\labelenumi{\arabic{enumi}.}
\setcounter{enumi}{9}
\tightlist
\item
  \begin{enumerate}
  \def\labelenumii{(\alph{enumii})}
  \tightlist
  \item
    Let \(\mathbf{A}'\) be the polynomial ring \(\mathbf{Z}[\mathbf{X},\mathbf{Y}]\) in two indeterminates and let \(\sigma\) be the automorphism of \(\mathbf{A}'\) leaving \(\mathbf{Y}\) invariant and such that \(\sigma (\mathbf{X}) = -\mathbf{X}\) ; \(\sigma\) and the identity automorphism form a group \(\mathcal{G}\) of automorphisms of \(\mathbf{A}'\) and the ring of invariants of \(\mathcal{G}\) is \(\mathbf{A} = \mathbf{Z}[\mathbf{X}^2,\mathbf{Y}]\) . Let \(\mathfrak{p}'\) be the prime ideal \(\mathrm{A}'(\mathrm{Y} - 2\mathrm{X})\) , \(\mathfrak{p} = \mathrm{A}\cap \mathfrak{p}'\) ; the decomposition group \(\mathcal{G}^{\mathbb{Z}}(\mathfrak{p}')\) is reduced to the identity element. Show that \(\mathrm{A / p}\) and \(\mathrm{A}' / \mathfrak{p}'\) are not isomorphic (consider the tensor products of these rings with \(\mathbf{Z} / 2\mathbf{Z}\) ).
  \end{enumerate}
\end{enumerate}

\begin{enumerate}
\def\labelenumi{(\alph{enumi})}
\setcounter{enumi}{1}
\tightlist
\item
  Let K be a field of characteristic \(\neq2\) , A' the polynomial ring K{[}X, Y{]} and G the automorphism group of A' consisting of the identity and the K{[}Y{]}-automorphism \(\sigma\) of A' such that \(\sigma(\mathbf{X}) = -\mathbf{X}\) ; the ring of invariants A of G is \(K[X^{2}, Y]\) . Let \(p'\) be the prime ideal \(\mathbf{A}'(\mathbf{X}^{3} - \mathbf{Y})\) of A' and \(p = A \cap p'\) ; the decomposition group \(\mathcal{G}^{\mathrm{Z}}(p')\) is reduced to the identity element but A/p and \(A'/p'\) are not isomorphic.
\end{enumerate}

¶ 11. (a) Let A be a local integrally closed domain, m its maximal ideal, K its field of fractions, K' a finite Galois extension of K, A' the integral closure of A in K', m' a maximal ideal of A' lying above m, B = \(\mathbf{A}^{\mathbb{Z}}(\mathfrak{m}')\) its decomposition ring and n = m' ∩ B. For every integer k \textgreater{} 0, B = A + n\^{}k (argue as in Proposition 4). Moreover there exists a primitive element z of \(\mathbf{K}^{\mathbf{Z}}\) (field of fractions of B) such that z ∈ n; deduce that n ∩ A{[}z{]} = mA{[}z{]}.

\begin{enumerate}
\def\labelenumi{(\alph{enumi})}
\setcounter{enumi}{1}
\item
  Show that, for every integer \(k\) , \(\mathfrak{m}^k\mathrm{B}_{\mathfrak{n}} \cap \mathrm{A} = \mathfrak{m}^k\) . (If \(x \in \mathfrak{m}^k\mathrm{B}_{\mathfrak{n}} \cap \mathrm{A}\) , note that there exists \(t \in \mathrm{A}[z]\) (where \(z\) is defined in (a)) such that \(t \notin \mathfrak{n}\) and \(tx \in \mathfrak{m}^k\mathrm{A}[z]\) ; writing \(t = t_0 + t_1z + \dots + t_{r-1}z^{r-1}\) (where \(r\) is the degree \([\mathbf{K}^{\mathbf{Z}}: \mathbf{K}]\) ), deduce that \(t_0x \in \mathfrak{m}^k\) and \(t_0 \notin \mathfrak{m}\) .)
\item
  Suppose now that \(\mathbf{K}'\) is any Galois extension of \(\mathbf{K}\) . Keeping the same notation as above, show that \(\mathbf{B} = \mathbf{A} + \mathfrak{n}^k\) and \(\mathfrak{m}^k\mathrm{B}_{\mathfrak{n}} \cap \mathbf{A} = \mathfrak{m}^k\) still (note that every element of \(\mathbf{A}^{\mathbb{Z}}(\mathfrak{m}')\) is contained in the decomposition ring of \(\mathfrak{m}' \cap \mathbf{C}\) , where \(\mathbf{C}\) is the integral closure of \(\mathbf{A}\) in a finite Galois sub-extension of \(\mathbf{K}'\) ).
\item
  Under the hypotheses of (c), show that, if A is Noetherian, so is \(B_{n}\) . (Observe that the Hausdorff completion of \(B_{n}\) with respect to the m-adic topology is identified with the completion \(\hat{A}\) of A by (c); if \(\phi: B_{n} \rightarrow \hat{A}\) is the
\end{enumerate}

canonical homomorphism, show that, for every ideal \(a\) of \(B_n\) , \(\phi(a\hat{A}) = a\) , using the fact that in \(\hat{A}\) every ideal is finitely generated and hence that \(a\hat{A}\) is generated by a finite number of elements of \(a\) and thus reduce it to the case where \(K'\) is a finite extension of \(K\) .)

\begin{enumerate}
\def\labelenumi{\arabic{enumi}.}
\setcounter{enumi}{11}
\tightlist
\item
  Let K be a field, \(K'\) a finite Galois extension of K, C the polynomial ring \(K'[X, Y]\) , B the quotient ring C/CXY and x, y the canonical images of X, Y in B. Let \(A'\) be the subring \(K[x] + yK'[y]\) of B. The Galois group G of \(K'\) over K operates faithfully on \(A'\) , the ring of invariants A of G in \(A'\) being \(K[x, y]\) . Let S be the set of powers of x in \(A'\) . Show that \(S^{-1}A' = S^{-1}A\) and that G does not operate faithfully on \(S^{-1}A'\) .
\end{enumerate}

¶ 13. (a) Let K be a field of characteristic 0 and n the ideal of the polynomial ring K{[}X, Y, Z{]} generated by \(Y^{2} - X^{2} - X^{3}\) . Show that n is prime and that the integral domain A = K{[}X, Y, Z{]}/n is not integrally closed: if x, y, z are the images of X, Y, Z in A, the element t = y/x of the field of fractions L of A is integral over A but does not belong to A. Let \(A' \supset A\) be the subring K{[}x, t, z{]} of L, which is integral over A. Consider in A the prime ideal p generated by xz - y and \(z^{2} - 1 - x\) and the maximal ideal \(q \supset p\) generated by x, y, z - 1. Consider on the other hand in A' the maximal ideal \(q'\) generated by x, \(t + 1\) and z - 1; show that \(q'\) lies above q but that there exists no prime ideal \(p' \subset q'\) of A' lying above p.~(Consider a homomorphism from \(A'/pA'\) to an extension field of K and show that the image of \(q'\) under such a homomorphism is necessarily 0.)

\begin{enumerate}
\def\labelenumi{(\alph{enumi})}
\setcounter{enumi}{1}
\item
  Let \(\mathbf{A}'\) be the quotient ring of the polynomial ring \(\mathbf{Z}[\mathbf{X}]\) by the ideal \(n\) generated by 2X and \(\mathbf{X}^2 -\mathbf{X}\) ; let \(x\) be the image of \(\mathbf{X}\) in \(\mathbf{A}'\) ; \(\mathbf{A}'\) is integral over \(\mathbf{Z}\) but the number 2 is a divisor of 0 in \(\mathbf{A}'\) . Let \(q'\) be the prime ideal of \(\mathbf{A}'\) generated by 2 and \(x - 1\) ; then \(q' \cap \mathbf{Z} = 2\mathbf{Z}\) ; if \(p\) is the prime ideal (0) in \(\mathbf{Z}\) , show that there exists no prime ideal \(p' \subset q'\) contained in \(\mathbf{A}'\) which lies above \(p\) .
\item
  Let K be a field, n the ideal of the polynomial ring K{[}X, Y, Z{]} generated by \(X^{2}\) , XY, XZ, YZ - Y and \(Z^{2} - Z\) , A' the quotient ring K{[}X, Y, Z{]}/n, x, y, z the images of X, Y, Z in A' and A the subring K{[}x, y{]} of A'; A is integrally closed in its total ring of fractions and A' is integral over A. Consider in A the prime ideal p generated by x and the maximal ideal q generated by x and y. Consider on the other hand in A' the maximal ideal \(q'\) generated by x, y, z; show that \(q'\) lies above q but that there exists no prime ideal \(p' \subset q'\) lying above p (same method as in (a)).
\end{enumerate}

* Interpret the examples in (a) and (c) geometrically. *

¶ 14. (a) Let A be a subring of the field K and x an element \(\neq0\) of K. Show that, if x is not integral over A, there exists a maximal ideal of \(A[x^{-1}]\) which contains \(x^{-1}\) and that, for every maximal ideal m of \(A[x^{-1}]\) containing x, the ideal \(A\cap m\) is maximal in A (observe that x does not belong to the ring \(A[x^{-1}]\) ).

\begin{enumerate}
\def\labelenumi{(\alph{enumi})}
\setcounter{enumi}{1}
\item
  Let A be a local integral domain, m its maximal ideal, K a field containing A and k the residue field of A. Let \(x \in K\) be such that \(x \notin A\) and \(x^{-1} \notin A\) and suppose that A is integrally closed in the ring A{[}x{]}. Show that there exists an A-algebra homomorphism \(A[x] \to k[X]\) which extends the canonical homomorphism \(A \to k\) and which maps x to X. (It is necessary to prove that, if \(\sum_{i=0}^{n}a_{i}x^{i}=0\) , where \(a_{i}\in A\) for all i, then all the \(a_{i}\) belong to m. Otherwise there would be a relation \(a_{n}x^{n-p}+\cdots+a_{p}=-(a_{p+1}x^{-1}+\cdots+a_{0}x^{-p})\) , where \(a_{p}\) is invertible in A; using Exercise 5 of §1, show that the common value b of the two sides belongs to A and, using (a), show that necessarily \(b\in m\) ; conclude that \(x^{-1}\) would be integral over A, which is absurd.)
\item
  Under the hypotheses of (b), show that the ideal mA{[}x{]} of A{[}x{]} is prime and not maximal.
\end{enumerate}

\begin{enumerate}
\def\labelenumi{\arabic{enumi}.}
\setcounter{enumi}{14}
\item
  Let A be a ring and p a prime ideal of A. Show that, in the polynomial ring B = A{[}X{]}, the ideal pB = q is prime; the maximal ideal of the local ring \(B_{q}\) is equal to \(pB_{q}\) ; if k is the field of fractions of A/p, \(B_{q}/pB_{q}\) is isomorphic to the field of rational fractions \(k(\mathbf{X})\) . If A is integrally closed, so is \(B_{q}\) .
\item
  Let A be an integral domain, K its field of fractions, \(K'\) a finite extension of K of degree n and \(A'\) a subring of \(K'\) containing A, with \(K'\) as field of fractions and integral over A. Let m be a maximal ideal of A and B the local ring \(A[X]_{mA[X]}\) (Exercise 15).
\end{enumerate}

\begin{enumerate}
\def\labelenumi{(\alph{enumi})}
\item
  The subring \(\mathbf{B}' = \mathbf{B}[\mathbf{A}']\) of \(\mathbf{K}'(\mathbf{X})\) is integral over \(\mathbf{B}\) ; its field of fractions is \(\mathbf{K}'(\mathbf{X})\) and \([\mathbf{K}'(\mathbf{X}) : \mathbf{K}(\mathbf{X})] = n\) .
\item
  For every maximal ideal \(\mathfrak{m}'\) of \(\mathbf{A}'\) containing \(\mathfrak{m}\) , \(\mathfrak{n}' = \mathfrak{m}'\mathbf{B}'\) is a maximal ideal of \(\mathbf{B}'\) ; if we write \(\mathfrak{n} = \mathfrak{m}\mathbf{B}\) , then
\end{enumerate}

\[
[ (\mathbf {A} ^ {\prime} / \mathfrak {m} ^ {\prime}): (\mathbf {A} / \mathfrak {m}) ] = [ (\mathbf {B} ^ {\prime} / \mathfrak {n} ^ {\prime}): (\mathbf {B} / \mathfrak {n}) ]
\]

and

\[
[ (\mathrm{A} ^ {\prime} / \mathfrak {m} ^ {\prime}): (\mathrm{A} / \mathfrak {m}) ] _ {s} = [ (\mathrm{B} ^ {\prime} / \mathfrak {n} ^ {\prime}): (\mathrm{B} / \mathfrak {n}) ] _ {s}.
\]

(Observe that \((\mathbf{A}' / \mathfrak{m}') \otimes_{\mathbf{A}} \mathbf{B}\) is isomorphic to the field \((\mathbf{A}' / \mathfrak{m}')(\mathbf{X})\) and deduce that \(\mathbf{B}' / \mathfrak{n}'\) is isomorphic to \((\mathbf{A}' / \mathfrak{m}')(\mathbf{X}).\) )

\begin{enumerate}
\def\labelenumi{(\alph{enumi})}
\setcounter{enumi}{2}
\tightlist
\item
  The mapping \(m' \mapsto m'B'\) is the bijection of the set of maximal ideals of \(A'\) containing m onto the set of maximal ideals of \(B'\) ; the inverse bijection is \(n' \mapsto n' \cap A'\) .
\end{enumerate}

¶ 17. (a) Let K be an infinite field, E a finite separable algebraic extension of K and g a polynomial of K{[}X{]}. Show that there exists an element \(x \in E\) such that \(\mathrm{E} = \mathrm{K}(x)\) and the minimal polynomial of x over K is relatively prime to g (argue as in Algebra, Chapter V, § 7, no. 7, Proposition 12).

\begin{enumerate}
\def\labelenumi{(\alph{enumi})}
\setcounter{enumi}{1}
\item
  Let \(k\) be a field, A a primary \(k\) -algebra (Chapter IV, § 2, Exercise 32) and \(m\) its unique prime ideal. If \(A / m\) is a transcendental extension of \(k\) , there exists in A elements which are not algebraic over \(k\) . If \(A / m\) is algebraic over \(k\) and \(d\) is an integer at least equal to the separable factor of the degree of \(A / m\) over \(k\) (and hence any integer if this separable factor is infinite), show that there exists in A elements whose minimal polynomial over \(k\) is of degree \(\geqslant d\) . If further \(A / m\) is not separable over \(k\) or if \(m \neq 0\) , there exists in A an element whose minimal polynomial over \(k\) is of degree \(>d\) . (If \(A / m\) is not separable over \(k\) , use Algebra, Chapter V, § 8, Exercise 5. If \(A / m\) is separable and of degree \(d\) over \(k\) and \(x \in A\) is such that \(\xi \in A / m\) , the class of \(x\) , is a primitive element of \(A / m\) , whose minimal polynomial is \(f \in k[X]\) , observe that \(f'(\xi) \neq 0\) and deduce that, if \(y \neq 0\) is an element of \(m\) , then \(f(x + y) \neq 0\) .)
\item
  Let \(k\) be an infinite field and \(A\) a \(k\) -algebra the direct composition of \(r\) primary algebras \(A_{i} (1 \leqslant i \leqslant r)\) ; let \(m_{i}\) be the maximal ideal of \(A_{i}\) and \(d_{i}\) an integer at least equal to the separable factor of the degree of \(A_{i} / m_{i}\) over \(k\) if \(A_{i}\) is algebraic and this factor is finite, or an arbitrary integer otherwise. Show that there exists in \(A\) an element which is not algebraic over \(k\) or
\end{enumerate}

whose minimal polynomial over \(k\) is of degree \(\geqslant \sum_{i=1}^{r} d_i\) ; if further one of the \(m_i\) is \(\neq 0\) or if one of the fields \(A_i / m_i\) is not a separable algebraic extension of \(k\) , there exists in \(A\) an element which is not algebraic over \(k\) or whose minimal

polynomial over \(k\) is of degree \(> \sum_{i=1}^{r} d_i\) (use (a) and (b)).

¶ 18. Let A be a local integrally closed domain, K its field of fractions, m its maximal ideal and \(k = \mathbf{A} / \mathfrak{m}\) its residue field. Let \(\mathbf{K}'\) be a finite algebraic extension of \(\mathbf{K}\) of degree \(n\) , \(\mathbf{A}'\) a subring of \(\mathbf{K}'\) containing \(\mathbf{A}\) , integral over \(\mathbf{A}\) and with \(\mathbf{K}'\) as field of fractions and \(\mathfrak{m}_i\) ( \(1 \leqslant i \leqslant r\) ) the maximal ideals of \(\mathbf{A}'\) . An ideal \(\mathfrak{m}\) is called unramified in \(\mathbf{A}'\) if there exists a basis \((w_i)_{1 \leqslant i \leqslant n}\) of \(\mathbf{K}'\) over \(\mathbf{K}\) consisting of elements of \(\mathbf{A}'\) and whose discriminant \(D_{\mathbf{K}' / \mathbf{K}}(w_1, \ldots, w_n)\) (Algebra, Chapter IX, §2) belongs to \(\mathbf{A} - \mathfrak{m}\) .

\begin{enumerate}
\def\labelenumi{(\alph{enumi})}
\item
  Show that, if \(\mathfrak{m}\) is unramified in \(\mathbf{A}'\) , \(\mathbf{K}'\) is a separable extension of \(\mathbf{K}\) , \(\mathbf{A}'\) is the integral closure of \(\mathbf{A}\) in \(\mathbf{K}'\) and is a free \(\mathbf{A}\) -module every basis of which over \(\mathbf{A}\) is a basis of \(\mathbf{K}'\) over \(\mathbf{K}\) ; \(\mathfrak{mA}'\) is the Jacobson radical of \(\mathbf{A}'\) and \(\mathbf{A}' / \mathfrak{mA}'\) is a semi-simple \(k\) -algebra of rank \(n\) over \(k\) which is separable over \(k\) and the direct composition of the fields \(\mathbf{A}' / \mathfrak{m}_i'\) .
\item
  Let \(d_{i}\) be the separable factor of the degree of \(A'/m_{i}'\) over \(A/m\) . Show that, if \(k\) is infinite, \(\sum_{i=1}^{r} d_{i} \leqslant n\) (use Exercise 17 (c)); moreover, the following conditions are equivalent:
\end{enumerate}

(α) m is unramified in A'.

\((\beta)\) \(\mathbf{A}'\) is a finitely generated A-module and \(\mathbf{A}' / \mathfrak{mA}'\) is a semi-simple \(k\) -algebra and separable over \(k\) .

\((\gamma)\sum_{i=1}^{r}d_{i}\leqslant n.\)

(To see that \((\beta)\) implies \((\gamma)\) , note that there is a system of generators of the A-module \(A'\) whose classes mod. \(mA'\) form a basis of \(A'/mA'\) over \(k\) , using Chapter II, § 3, no. 2, Corollary 2 to Proposition 4; deduce that

\[
[ \mathbf {K} ^ {\prime}: \mathbf {K} ] \leqslant [ (\mathbf {A} ^ {\prime} / \mathsf {m A} ^ {\prime}): k ].
\]

To see that \((\gamma)\) implies \((\alpha)\) , use Exercise 17 (c) to show that there exists a basis of \(\mathbf{A}' / \mathfrak{mA}'\) over \(k\) consisting of the powers \(\xi^i\) of an element \(\xi\) ( \(0 \leqslant i \leqslant n - 1\) ); if \(x \in \mathbf{A}'\) is an element of the class \(\xi\) , deduce that \(\mathrm{D}_{\mathrm{K}' / \mathrm{K}}(1, x, \ldots, x^{n-1})\) belongs to \(\mathbf{A} - \mathfrak{m}\) .) Give an example where \(\mathbf{K}'\) is separable over \(\mathbf{K}, \mathbf{A}' / \mathfrak{mA}' = \mathbf{A} / \mathfrak{m}\) and \(\mathbf{A}'\) is not a finitely generated A-module (cf.~Exercise 9 (c)).

\begin{enumerate}
\def\labelenumi{(\alph{enumi})}
\setcounter{enumi}{2}
\tightlist
\item
  Extend the results of (b) to the case where \(k\) is finite (in the notation of Exercise 16, show that, if \(n\) is unramified in \(B'\) , then \(m\) is unramified in \(A'\) ; taking account of the fact that \(B / n\) is infinite, use Exercise 16 (c), thus obtaining an element \(z \in B'\) such that
\end{enumerate}

\[
\mathrm{D} _ {\mathbf {K} ^ {\prime} (\mathbf {X}) / \mathbf {K} (\mathbf {X})} (1, z, \dots , z ^ {n - 1}) \in \mathrm{B} - n;
\]

show that \(z\) may be assumed to be a polynomial in \(\mathbf{A}'[\mathbf{X}]\) and, by expressing the above discriminant as a polynomial in \(\mathbf{X}\) , obtain the existence of a system of \(n\) elements \(w_{i}\) ( \(1 \leqslant i \leqslant n\) ) of \(\mathbf{A}'\) such that \(\mathrm{D}_{\mathrm{K}'/\mathrm{K}}(w_1, \ldots, w_n) \in \mathbf{A} - \mathfrak{m}\) ).

¶ 19. Let A be an integrally closed domain, K its field of fractions, K' a finite algebraic extension of K of degree n and A' a subring of K' containing A, integral over A and with K' as field of fractions; if p is a prime ideal of A, let \(\mathbf{A}_{\mathfrak{p}}^{\prime}\) denote the ring of fractions of \(\mathbf{A}^{\prime}\) whose denominators are in \(\mathbf{A}-\mathfrak{p},\mathfrak{p}\) is called unramified in \(\mathbf{A}^{\prime}\) if \(\mathfrak{p}\mathbf{A}_{\mathfrak{p}}\) is unramified in \(\mathbf{A}_{\mathfrak{p}}^{\prime}\) .

\begin{enumerate}
\def\labelenumi{(\alph{enumi})}
\tightlist
\item
  Let \(\mathfrak{p}_i^{\prime}(1\leqslant i\leqslant r)\) be the prime ideals of \(\mathbf{A}'\) lying above \(\mathfrak{p}\) . Let \(d_{i}\) be the separable factor of the degree of the field of fractions of \(\mathbf{A}' / \mathfrak{p}_i'\) over the field of
\end{enumerate}

fractions of A/p; show that \(\sum_{i=1}^{r} d_i \leqslant n\) . (Reduce it to Exercise 18.)

\begin{enumerate}
\def\labelenumi{(\alph{enumi})}
\setcounter{enumi}{1}
\tightlist
\item
  The following conditions are equivalent:
\end{enumerate}

(α) The ideal p is unramified in A'.

(β) A' contains a basis of K' over K whose discriminant belongs to A - p.

\((\gamma)\sum_{i=1}^{r}d_{i}=n.\)

In particular, if \(r = n\) , \(\mathfrak{p}\) is unramified in \(\mathbf{A}'\) .

Give an example where \(p\) is unramified in \(A'\) but \(A'\) is not a finitely generated \(A\) -module (cf.~Exercise 9 (a)).

\begin{enumerate}
\def\labelenumi{(\alph{enumi})}
\setcounter{enumi}{2}
\item
  A is called unramified in A' (or in K') if every prime ideal of A is unramified in A'. Show that A' is then a faithfully flat A-module (cf.~Chapter II, § 3, no. 4, Corollary to Proposition 15). Let K'' be a finite algebraic extension of K' of degree n'; let A' be the integral closure of A in K' and A'' the integral closure of A in K''; for A to be unramified in A'', it is necessary and sufficient that A be unramified in A' and A' be unramified in A''.
\item
  Let \(k\) be a field, \(A\) an integrally closed integral \(k\) -algebra and \(K\) its field of fractions. Let \(k'\) be a finite separable algebraic extension of \(k\) of degree \(n\) over \(k\) . Show that, if \(k' \otimes_k K\) is a field, \(A\) is unramified in \(k' \otimes_k K\) (consider a primitive element \(x\) of \(k'\) over \(k\) and the basis of \(k' \otimes_k K\) consisting of the \(x^i\) for \(0 \leqslant i \leqslant n - 1\) ).
\end{enumerate}

\begin{enumerate}
\def\labelenumi{\arabic{enumi}.}
\setcounter{enumi}{19}
\item
  Let A be the integral domain defined in Chapter III, § 3, Exercise 15 (b), C the local ring \(\mathbf{A}_{\mathfrak{m}}\) , \(\mathfrak{n} = \mathfrak{mA}_{\mathfrak{m}}\) the maximal ideal of C and L the field of fractions of A and C. If \(x, y\) are the canonical images of X and Y in C, show that \(t = y / x\) does not belong to C but that \(t^2 \in \mathbf{C}\) , so that C is not integrally closed; moreover \(\mathbf{C}' = \mathbf{C}[t]\) is the integral closure of C and is isomorphic to \(\mathbf{S}^{-1}\mathbf{K}[\mathbf{T}]\) (T an indeterminate), where S is the multiplicative subset of K{[}T{]} consisting of the polynomials which are not divisible by T - 1 nor by T + 1. Show that \(\mathbf{C}'\) is a finitely generated C-module and that the (maximal) ideals of \(\mathbf{C}'\) lying above \(\mathfrak{n}\) are the principal ideals \(\mathfrak{q}_1 = (t - 1)\mathbf{C}'\) and \(\mathfrak{q}_2 = (t + 1)\mathbf{C}'\) ; but the local rings \(\mathbf{C}_{\mathfrak{q}_1}'\) and \(\mathbf{C}_{\mathfrak{q}_2}'\) are not C-algebras integral over C.
\item
  Let A be an integrally closed domain containing the field Q, K its field of fractions, L an algebraic extension of K and B the integral closure of A in L. Show that, for every ideal a of A, \(A \cap aB = a\) .
\item
  Let A be a Noetherian integral domain, K its field of fractions and A' a subring of K containing A. Suppose that A is integrally closed in A' and that for every prime ideal p of A there exists a prime ideal \(p'\) of \(A'\) lying above p.~Argue by reductio ad absurdum: let \(x \in A' \cap C A\) and let \(a = A \cap (x^{-1} A)\) , which is an ideal distinct from A. Let p be a prime ideal associated with A/a; the transporter b = a: p is then distinct from a (Chapter IV, §1, no. 4, Proposition 9) and there therefore exists \(y \in b\) such that \(xy \notin A\) and
\end{enumerate}

\[
x y p \subset p A ^ {\prime} \cap A = p;
\]

deduce a contradiction.)

